\paragraph{Comparison with bilinear agreement bounds}
\label{sec:profile-comparison}
Cohen and Moshkovitz establish
$\operatorname{AR}(D)\ge(1-\log_q2)\operatorname{GR}(D)$ and
$\operatorname{AR}(D)=-\log_q\Pr[D=0]$ for uniform bilinear inputs
\cite[Proposition~4.1 and Fact~4.5]{cohen2022structure}. Rectangular multiplication has geometric
rank $g=nr-\lfloor r^2/4\rfloor$~\cite{kopparty2023geometric}
and slice rank $nr$ even over the algebraic closure, by its
semistability~\cite[Example~24 and Lemma~23]{christandl2023weighted}.\footnote{If a slice decomposition uses dimensions $a_i$ in modes of
dimensions $m_i$ with $\sum_i a_i<\min_i m_i$, assign weights
$1-a_i/m_i$ to its subspaces and $-a_i/m_i$ to their complements.
After clearing denominators, each factor has total weight zero, while
every supported tensor term has positive weight
$\ge1-\sum_i a_i/m_i$. This contradicts semistability. The reverse
bound follows by slicing a smallest mode.}
For $D=T-F$, take slice rank over the algebraic closure.
Subadditivity gives $\operatorname{SR}(D)\ge nr-w$ and
$\operatorname{GR}(D)\ge g-w$. Cohen--Moshkovitz's proof gives
$\operatorname{SR}(D)\le2\operatorname{GR}(D)$ over that closure
\cite[proof of Theorem~3.1]{cohen2022structure}; no descent is needed.
Since geometric rank is integral, their finite-field analytic-rank
inequality yields
\begin{equation}\label{eq:cm-field-comparator}
 \Pr[F=T]\le (2/q)^{\max\{\lceil(nr-w)/2\rceil,\,g-w,\,0\}},\quad q>2.
\end{equation}
An elementary comparison also uses the exact zero probability:
a uniform-input bilinear map's largest output atom occurs at zero
\cite[proof of Proposition~5.1]{cohen2022structure}. The fixed range of
$F$ has at most $q^w$ points, giving
$B_{\rm range}(w)=\min\{1,q^w\mu(0)\}$. Concavity of $z\mapsto\min\{1,z\}$ gives
\begin{equation}\label{eq:profile-vs-range}
 \mu(w)=\E_s\min\{1,q^{w-ns}\}\le B_{\rm range}(w).
\end{equation}
This is strict for every integer $0<w<nr$: the positive rank-zero
mass makes $\mu(w)<q^w\mu(0)$, and positive full-rank mass makes
$\mu(w)<1$. At $w=0$ the two range bounds coincide exactly.
At $n=16,r=4,q=17,w=56$, $\mu(w)=1.434\times10^{-10}$,
$B_{\rm range}(w)=1.426\times10^{-5}$, and
\eqref{eq:cm-field-comparator} gives $1.916\times10^{-4}$.
The rank-conditioned argument also covers $q=2$, where the geometric
comparison's coefficient vanishes. Below a certified tensor-rank
threshold, all three bounds can also use the common cap $2/q-1/q^2$;
checking then applies $p\mapsto p+(1-p)q^{-t}$.
